\section{Introduction}
\label{sec:introduction}

\subsection{The Ensemble Bottleneck}

Following \textit{Rucho v.\ Common Cause}~\citep{scotus2019rucho}, partisan
gerrymandering claims have migrated from federal to state courts, and Markov
chain ensemble analysis has become the principal expert-witness tool for
showing that an enacted map is a statistical outlier.
The dominant workflow runs thousands of algorithmically drawn plans using
the \recom{} chain~\citep{deford2021recombination}, then places the contested
map within the resulting distribution.
A plan at the 99th percentile of Republican seat share is characterized as
strong evidence of partisan intent; at the 50th percentile, it is
indistinguishable from a random draw.

The limiting factor is throughput.
The primary open-source implementation, \gc{}~\citep{gerrychain2021},
is written in Python.
For North Carolina (2{,}672 census tracts, $k=14$ districts), \gc{}
achieves approximately 21 steps per second.
A credible 10{,}000-step ensemble therefore requires roughly 8 minutes;
a 50{,}000-step ensemble requires 40 minutes.
At this speed, ensemble analysis is necessarily a separate, overnight
computation---a post-hoc audit distinct from the redistricting pipeline
itself.

This throughput ceiling creates two practical problems.
First, ensemble analysis cannot serve as a \emph{real-time} feasibility
check during map-drawing.
A redistricting practitioner cannot interactively explore whether a
proposed boundary change moves the plan toward or away from the ensemble
median.
Second, the computational cost discourages the thorough convergence
diagnostics---multiple parallel chains, extended runs---that would give
courts and litigants justified confidence in the ensemble's stationarity
(see~\citealt{dellaLibera2026g4}).

\subsection{The Rust Solution}

We present \re{}, a Rust implementation of the \recom{} Markov chain
built around Wilson's uniform random spanning tree
algorithm~\citep{wilson1996generating}.
The design goal is to collapse the throughput gap sufficiently that
a 10{,}000-step ensemble for any congressional-scale state completes in
under one second, enabling ensemble analysis to run as an integrated step
within \texttt{redist build} rather than as a separate overnight job.

The key algorithmic innovation is the use of Wilson's loop-erased random
walk~\citep{lawler1980loop} to generate uniform random spanning trees.
Wilson's algorithm runs in $O(\tau_{\mathrm{cover}})$ expected time, where
$\tau_{\mathrm{cover}}$ is the cover time of the underlying random walk.
For planar graphs---the natural setting for census-tract adjacency
graphs---this is $O(|V|\log|V|)$~\citep{aldous1991random}.
Each \recom{} step uses the spanning tree of the merged pair region
(approximately $2n/k$ nodes, where $n$ is the total tract count and $k$
is the number of districts), so Wilson's cost per step is
$O((n/k)\log(n/k))$.
Rust's compiled, zero-overhead arithmetic exploits this theoretical
efficiency to an extent that interpreted Python cannot.

\subsection{Pipeline Integration}

The intended deployment is as the \texttt{redist ensemble} subcommand,
integrated into the same binary as \texttt{redist state},
\texttt{redist analyze}, and \texttt{redist map}.
A typical invocation:

\begin{verbatim}
  redist ensemble --state NC --year 2020 --steps 10000 \
                  --chains 8 --seed 42 --out nc_ensemble.json
\end{verbatim}

At 50{,}000 steps/sec, this completes in approximately 0.2 seconds for
North Carolina.
The output JSON carries $\rhat$, $\ess$, and Hamming autocorrelation
diagnostics, enabling immediate convergence assessment without a separate
analysis pass.

\subsection{Paper Organization}

Section~\ref{sec:related} reviews \recom{}, Wilson's spanning tree
algorithm, and \gc{}.
Section~\ref{sec:algorithm} gives the formal algorithmic description.
Section~\ref{sec:implementation} describes the Rust implementation
architecture.
Section~\ref{sec:performance} presents the performance comparison.
Section~\ref{sec:diagnostics} describes the convergence diagnostics.
Section~\ref{sec:conclusion} concludes.
